\chapter{Results, Evaluation and Discussion}
\label{chapter5}

This chapter is intentionally structured in advance of the final authoritative matched-arm execution. Numerical content will be inserted after running the redesigned Stage B experiment and collecting the final output artifacts. Until then, the chapter records the reporting structure, the interpretation plan, and the evidence sources that will be used.


\section{Online Experiment Results}
\label{sec:results}

The final experiment will populate this chapter from the redesigned Stage B artifacts: the baseline and custom arm CSV files, the per-run metadata file, the exported node-capacity JSON, the decision-artifact directory, and the final \texttt{stage\_b\_summary.json}. The reporting emphasis will remain on custom-minus-baseline deltas so that the reader can distinguish scheduler effect from cluster-wide noise.

\subsection{Aggregate Metrics}

This subsection will report the aggregate custom-minus-baseline deltas for:

\begin{itemize}
\item safe placement rate
\item anomalous placement rate
\item high-contention rate
\item scheduling latency
\item startup time
\item total time
\item placement fairness
\item capacity alignment
\item utility score
\end{itemize}

One summary table will be inserted here after the run. The accompanying text will interpret whether the custom arm provides a net placement-quality improvement and whether any latency cost is operationally meaningful.


\subsection{Per-Phase Breakdown}

This subsection will summarise the experiment by phase: normal, stress\_cpu, stress\_memory, and stress\_mixed. The primary purpose is to test whether the custom scheduler helps most under stressed operation, which is the condition it was designed for.

The post-run analysis will report phase-level deltas for timing, contention, placement safety, and capacity alignment. If phase-specific reversals occur, they will be discussed here rather than hidden inside aggregate averages.


\subsection{Per-Workload Breakdown}

This subsection will report the interaction between phase and workload type (CPU, memory, and mixed). The key question is whether the custom scheduler behaves differently across single-resource and mixed-resource pods once workload requests are injected directly into the live ranker.

The final version will include one table or figure summarising the phase$\times$workload matrix and a short paragraph identifying the strongest improvement bucket and the largest regression bucket, if any.


\subsection{Per-Decision Diagnostics}

The redesigned pipeline records decision-level metadata that was absent from the earlier Stage B path. This subsection will use those artifacts to report:

\begin{itemize}
\item expected-node match rate for the custom arm
\item stress-target avoidance or placement rates
\item prediction-source counts (trained model versus fallback)
\item mean decision score components, including projected load, anomaly risk, and capacity penalty
\end{itemize}

This material is important because it connects high-level outcome metrics back to the actual ranking behaviour of the custom arm.


\section{Hypothesis Evaluation}
\label{sec:hypothesis-eval}

This section will evaluate each hypothesis from Chapter~\ref{chapter3} against the final online evidence.

\subsection{H1: Placement Quality}

Placeholder for final verdict after the run. The interpretation will be based primarily on safe placement, anomalous placement, and high-contention deltas, with phase-specific evidence used to determine whether improvements are robust or restricted to particular stress conditions.


\subsection{H2: Scheduling Overhead}

Placeholder for final verdict after the run. The interpretation will compare scheduling latency cost against startup-time and total-time benefit, rather than treating latency in isolation.


\subsection{H3: Offline-to-Online Transfer}

Placeholder for final verdict after the run. The interpretation will consider whether the locked hybrid family produces positive online utility against the default scheduler and whether the decision-level diagnostics support that transfer claim.


\section{Discussion}
\label{sec:discussion}

\subsection{Effectiveness of the Hybrid Approach}

This subsection will interpret whether the custom scheduler's benefit arises mainly from anomaly avoidance, contention avoidance, workload-aware capacity normalisation, or a combination of the three. The redesigned decision artifacts should make this interpretation more defensible than in the earlier Stage B path.


\subsection{Scheduler Authority and Experimental Control}

The final discussion will explicitly address the methodological significance of the redesigned online path: explicit warm-up, counterbalanced arm order, authoritative rank-and-pin binding, rotating stress targets, and per-decision logging. These design changes should be discussed not as implementation ornamentation, but as the reason the final online claim is stronger than the earlier scheduler-profile comparison.


\subsection{Comparison with Related Work}

Placeholder for the final positioning discussion. This section will relate the final results to the anomaly-detection replication evidence in Appendix~\ref{chapter2} and to the broader scheduling literature reviewed in Chapter~\ref{chapter1}, while acknowledging that direct numerical comparison across platforms remains limited.


\subsection{Limitations}
\label{sec:limitations}

Some limitations are already known before the final run:

\begin{itemize}
\item the cluster is still modest relative to large production edge deployments
\item the live experiment relies on controlled stress injection rather than organic failure cascades
\item the offline replay surface does not reproduce every live execution refinement, especially the decision-time workload and capacity context
\item the predictor remains a lightweight linear family that may underfit complex mixed-resource regimes
\end{itemize}

After the run, this subsection will be extended with any result-specific limitations that become visible in the final artifacts.


\section{Conclusions}
\label{sec:conclusions}

Placeholder for the final experimental conclusion. The completed version will summarise the answer to RQ1--RQ3 using only the final matched-arm evidence and the locked offline support material described in Chapters~\ref{chapter3} and \ref{chapter4}.


\section{Future Work}
\label{sec:future-work}

The final version of this section will be updated after the run, but the likely directions remain:

\begin{itemize}
\item larger and more heterogeneous clusters
\item additional trials for tighter uncertainty estimates
\item more expressive predictors for mixed-resource regimes
\item more realistic anomaly and failure injection
\item deeper integration with native Kubernetes scheduling interfaces
\end{itemize}